% =====================================================================
%  SN3 (Supplementary Notes): what an executable skill is.
%  Fragment \input by supplement.tex under the Supplementary Notes part.
%  No preamble here.
% =====================================================================

\section{Executable skills}
\label{sec:supp-skills}

Because the term may be unfamiliar outside agent-based AI tooling, we briefly
define an executable \emph{skill}. A skill is a version-controlled package of
natural-language instructions and supporting resources, such as reference
files, templates, and worked examples. An AI assistant loads the package on
demand to perform a defined task in a repeatable way~\cite{anthropic_skills}.

A skill is the executable counterpart of a written protocol. A protocol tells
a human reader what a task requires; a skill expresses those requirements in a
form that an AI assistant can apply directly. In practice, a skill is a
folder containing a top-level instruction file, conventionally
\texttt{SKILL.md}, and the resources referenced by that file. The assistant
reviews the skills available to it, loads the one relevant to the task, and
follows its instructions. Its behavior is therefore governed by the
versioned package rather than improvised anew for each session.

The general mechanism is not specific to a particular model or vendor, and a
cross-vendor specification is emerging.\footnote{\url{https://agentskills.io}}
The two skills accompanying this work are implemented as Claude Agent Skills.

The annotation skill packages the annotation protocol, and the review skill
packages the review protocol (\cref{sec:supp-protocol-suite}). They allow an
AI assistant to apply the same rules that guide a human curator, under curator
supervision. Both skills are published with the protocols in the project
repository.\footnote{\url{https://github.com/ForomePlatform/genetic-evidence-model/tree/master/skills}}
Their purpose in this work is reproducibility and reuse, not autonomy. The
curator remains the authority for every recorded annotation, as described in
\cref{sec:supp-protocol-suite}.
